\section{Procesamiento de Señales digitales}
\subsection{Digitalización de una señal}
\begin{frame}{Señal digitalizada}
 \begin{textblock*}{\textwidth}(0.4cm,1.2cm)
  \begin{itemize}[<+->]\justifying
    \item El proceso de digitalización de una señal responde al siguiente modelo.
  \end{itemize}
  \only <2> {\centering \includegraphics[width=\textwidth]{digitalizacion-layer1.pdf} }
  \only <3> {\centering \includegraphics[width=\textwidth]{digitalizacion-layer2.pdf} }
  \only <4> {\centering \includegraphics[width=\textwidth]{digitalizacion-layer3.pdf} }
  \only <5> {\centering \includegraphics[width=\textwidth]{digitalizacion-layer4.pdf} }
  \only <6> {\centering \includegraphics[width=\textwidth]{digitalizacion-layer5.pdf} }
  \only <7> {\centering \includegraphics[width=\textwidth]{digitalizacion-layer6.pdf} }
  \only <8> {\centering \includegraphics[width=\textwidth]{digitalizacion-layer7.pdf} }
 \end{textblock*} 
 \begin{textblock*}{\textwidth}(0.4cm,6.8cm)
  \begin{itemize}[<+(1)->]\justifying
    \item La función \textit{v(t)}, es una tensión de entrada analógica.
    \item La función \textit{v'(t)}, es una función discreta, que solo cambia entre valores discretos (no contínua) a intervalos regulares, dados por la frecuencia de muestreo, y mantiene su valor durante el intervalo.
  \end{itemize} 
 \end{textblock*} 
\end{frame}

\begin{frame}{Señal digitalizada}
 \begin{textblock*}{\textwidth}(0.4cm,1.2cm)
  \begin{itemize}[<+->]\justifying
   \item El circuito de Muestreo y Retención sostiene durante todo el intervalo de muestreo el valor de la señal de entrada al inicio de dicho intervalo, ignorando los demás valores.
   \item De este modo actúa sobre la variable independiente de la señal de entrada, en nuestro caso el tiempo, convirtiéndola de continua a discreta, ya que define un conjunto finito de valores de tiempo que resultan de interés.
   \item El proceso de Cuantificación (o Cuantización), transforma los valores obtenidos de la variable dependiente aproximándolos al valor mas cercano de un conjunto finito de $2^{n}$ números.
   \item A continuación mas detalles\dots
  \end{itemize} 
 \end{textblock*} 
\end{frame}

\begin{frame}{1er. Fase: Muestreo y retención}
 \begin{textblock*}{0.5\textwidth}(0.3cm,1.3cm)
  \centering
  \includegraphics[width=\textwidth]{SandH.pdf}\\ \vspace{-0.3cm}
  \only <1> {\includegraphics[width=\textwidth]{SampleAndHold-layer1.pdf} }
  \only <2-> {\includegraphics[width=\textwidth]{SampleAndHold-layer2.pdf} }
 \end{textblock*}
 \begin{textblock*}{0.5\textwidth}(0.3cm+0.5\textwidth,1.0cm)
  \begin{itemize}[<+->]\justifying
   \item Se toma un valor instantáneo de la señal y se ``retiene'' el valor de tensión en una capacidad, con un circuito resistivo de descarga de muy alta resistencia.
   \item {\textit{\color{red}Comparando la señal saliente del circuito de Muestreo y Retención, con la señal original (en punteado suave), observar que el error aumenta en los cambios abruptos de señal.}}
   \item Otra forma de observar la relación entre la frecuencia de muestreo y la máxima frecuencia de la señal original
  \end{itemize}
 \end{textblock*}
\end{frame}

\begin{frame}{Cuantificación y Codificación}
 \begin{textblock*}{0.5\textwidth}(0.3cm,1.2cm)
  \centering \vspace{-0.3cm}\includegraphics[width=\textwidth]{Cuantif.pdf}\\
  \only <4-> { \vspace{-0.8cm}\includegraphics[width=\textwidth]{Cuantif-error.pdf} }
 \end{textblock*}
 \begin{textblock*}{0.5\textwidth}(0.3cm+0.5\textwidth,1.2cm)
  \begin{itemize}[<+->]\justifying
   \item Se llevan a cabo en un Conversor Analógico Digital.
   \item Se puede aproximar por redondeo o truncamiento
   \item Independientemente del método el error aumenta en las zonas donde la derivada de la señal es mas alta.
   \item Si tomamos la diferencia entre la señal muestreada y la cuantificada y operamos una diferencia se obtiene el error de Cuantificación o Cuantización)
  \end{itemize}
 \end{textblock*}
\end{frame}

\subsection{Arquitecturas de Procesamiento de una señal digital}
\begin{frame}{Procesador de Señales Digitales}
 \begin{textblock*}{\textwidth}(0.4cm,1.2cm)
  \begin{itemize}[<+(1)->]\justifying
   \item Es una CPU de propósito dedicado, diseñada para realizar cálculos y procesamiento de un único tipo de datos: secuencias de valores correspondientes a la codificación de las muestras de una señal de entrada.
   \item Su arquitectura está pensada para optimizar el procesamiento de datos que no son de gran tamaño (8, 16, 24, o a lo sumo 32 bis). Se trata de pixeles de una imagen, o de valores instantáneos de audio, o de señales médicas o de mapas térmicos, etc.
   \item La característica distintiva de este tipo de datos reside en sus algoritmos de cálculo: Normalmente se requiere procesar no solo el valor actual sino la combinación del valor actual con n valores anteriores en el tiempo, o vecinos (en el caso de una imagen lo que llamaremos N8)
  \end{itemize}
 \end{textblock*}
\end{frame}

\begin{frame}{Procesador de Señales Digitales}
 \begin{textblock*}{\textwidth}(0.4cm,1.2cm)
  \begin{itemize}[<+(1)->]\justifying
   \item En general una operación muy frecuente son los filtros de convolución, dados por una expresión del tipo:\\
$\displaystyle y[n] = \sum_{i=0}^{N} \; a_{n}\cdot x[n-i] = a_{0}\cdot x[n] + a_{1}\cdot x[n-1] + a_{2}\cdot x[n-2] + ...  + a_{N}\cdot x[n-N]$
   \item En general se deben resolver sumas de productos y acumular su resultado.
   \item Para optimizar se deberían leer y procesar varios datos en paralelo
   \item Las técnicas de paralelismo que se desarrollaron en estos procesadores se denominaron por tal razón \textbf{\textit{\color{blue}Data Level Paralelism}}.
  \end{itemize}
 \end{textblock*}
\end{frame}

\begin{frame}{Arquitectura de un Procesador de Señales Digitales}
 \begin{textblock*}{\textwidth}(0.4cm,1.2cm)
  \begin{itemize}[<+(1)->]\justifying
   \item En general el diseño de un Procesador de Señales Digitales concentra sus esfuerzos en resolver en paralelo:
   \begin{description}[<+(1)->]\justifying
    \item[El acceso a los operandos] Para este fin se implementan buses paralelos con una cantidad de líneas de datos superior al ancho de palabra de la CPU
    \item[El almacenamiento de resultados] Se resuelve mediante buses dedicados para manejar la salida de la ALU hacia memoria o registros (con las consideraciones que la concurrencia de accesos debe tener en cuenta en el hardware)
    \item[El procesamiento de la mayor cantidad de datos posible] Los primeros pasos se conocen como VLIW (Very Large Instruction Word), que derivó en el modelo de ejecución SIMD (Single Instruction Multiple Data)
   \end{description}
  \end{itemize}
 \end{textblock*}
\end{frame}
